\documentclass[10pt,letterpaper]{article}
\usepackage[T1]{fontenc}
\usepackage{lmodern}
\usepackage{microtype}
\usepackage{amsmath,amssymb,amsthm}
\usepackage{booktabs,tabularx}
\usepackage{xcolor}
\usepackage[margin=0.9in]{geometry}
\setlength{\emergencystretch}{1em}
\definecolor{linkblue}{RGB}{28,77,122}
\usepackage[colorlinks=true,allcolors=linkblue]{hyperref}

\newtheorem{proposition}{Proposition}

\hypersetup{
  pdftitle={A Size-Aware Directory of R*-Trees for Dynamic Rectangle Workloads},
  pdfsubject={Spatial Indexing}
}

\title{A Size-Aware Directory of R*-Trees for Dynamic Rectangle Workloads}
\author{}
\date{}

\begin{document}
\maketitle
\vspace{-2.5em}

\begin{abstract}
An R*-tree incurs substantial update overhead from overlap minimization and forced reinsertion. We study a fixed-depth quadtree directory that routes each planar rectangle to a single owning cell by its size and center, promotes busy cells from flat vectors to local R*-trees, and answers range and $k$-nearest-neighbor queries using conservative subtree bounding boxes. Across deterministic road-rectangle traces from 24 U.S.\ counties (capacity 32), the loose directory achieved a median replay speedup of $3.34\times$ (95\% paired bootstrap interval $3.11\text{--}3.46\times$) over an optimized single R*-tree with 30\% reinsertion ($2.55\times$, $2.44\text{--}2.68\times$ without reinsertion) and $1.25\times$ ($1.26\times$ without reinsertion) over a strict containing-cell directory. Gains ranged from $1.85\times$ on moving-hotspot traces to $6.48\times$ on insert-only traces with 30\% reinsertion. All 152 trace and reinsertion comparisons matched the archived query signatures and passed exhaustive verification on 316,717 primary queries per setting.
\end{abstract}

\section{Motivation and Design}

The R*-tree of Beckmann et al.~\cite{beckmann1990} guides node splitting and subtree selection using area, margin, and overlap criteria, supplemented by forced reinsertion. Maintaining that structure carries measurable update cost even for in-memory two-dimensional indexes. Our single-tree control already replaces quadratic split rescanning with prefix/suffix bound passes and applies exact early-out pruning to overlap accumulation once a candidate child cannot beat the incumbent. A further planar specialization preserved exact stock query and layout signatures while improving median replay speed by only $1.07\times$ with reinsertion. Achieving larger gains therefore required changing where insertions and splits occur.

We construct a fixed-depth ($D = 6$, i.e., $64 \times 64$ finest cells) quadtree directory whose padded bounding domain is fixed from the initial insertion prefix. Every rectangle is assigned a unique owner cell:
\begin{itemize}
\item \emph{Out-of-domain rectangles} (extending outside the padded root domain) are assigned to the root cell.
\item \emph{Strict mode} assigns an in-domain rectangle to the deepest cell that contains the entire rectangle.
\item \emph{Loose mode} assigns an in-domain rectangle to the cell containing its center at the deepest level $d \in \{0,\dots,D\}$ whose cell width and height are both at least those of the rectangle (so the rectangle fits inside the factor-$2$ expanded cell, in the spirit of loose quadtrees~\cite{samet2013}).
\end{itemize}
Each cell begins as a flat vector and promotes to a local R*-tree once an insertion would exceed the node capacity; moving an object removes the old rectangle from its owner and inserts the updated rectangle into its new owner.

\section{Conservative Subtree Bounds and Query Correctness}

Each directory node $u$ maintains a live descendant count $c(u)$ and a bounding rectangle $B(u)$. On insertion into cell $v$, every ancestor $u$ of $v$ increments $c(u)$ and expands $B(u)$ to cover the inserted rectangle (initializing $B(u)$ to the rectangle when $c(u)$ transitions from $0$ to $1$). On deletion, ancestors decrement $c(u)$ while leaving $B(u)$ unchanged.

\begin{proposition}[Soundness of conservative pruning]
At every state of the directory, $B(u)$ covers every live rectangle stored in the subtree rooted at $u$ whenever $c(u) > 0$. Consequently:
\begin{enumerate}
\item A range query $Q$ that skips any node $u$ with $c(u) = 0$ or $B(u) \cap Q = \emptyset$, and emits all live entries beneath any node $u$ with $B(u) \subseteq Q$ without per-entry intersection checks, returns the exact set of live rectangles intersecting $Q$.
\item A best-first $k$-nearest-neighbor search using squared Euclidean point-to-rectangle distance $\operatorname{dist}^2(q, B(u))$ as the priority key and terminating once $k$ candidates are held and the top queue key exceeds the $k$-th smallest observed distance returns a valid $k$-nearest-neighbor distance multiset.
\end{enumerate}
\end{proposition}

\begin{proof}
By induction on operations, every live rectangle in the subtree of $u$ was inserted at some step after the most recent transition of $c(u)$ from $0$ to $1$, at which point $B(u)$ was expanded to contain it; subsequent insertions only enlarge $B(u)$, and deletions do not shrink $B(u)$ while $c(u) > 0$. Therefore every live rectangle $R$ in the subtree of $u$ satisfies $R \subseteq B(u)$.

For range queries, $B(u) \cap Q = \emptyset \implies R \cap Q = \emptyset$, so skipped subtrees contain no matching entries; conversely, $B(u) \subseteq Q \implies R \subseteq Q \implies R \cap Q \neq \emptyset$, so all live descendants beneath a covered node match. For $k$-nearest-neighbor queries, $R \subseteq B(u)$ implies $\operatorname{dist}^2(q, B(u)) \leq \operatorname{dist}^2(q, R)$. Because the branch min-heap pops subtrees in nondecreasing order of $\operatorname{dist}^2(q, B(u))$, once $k$ live rectangles with squared distance at most $\tau$ have been found and the next heap key satisfies $\operatorname{dist}^2(q, B(u)) > \tau$, no unvisited live rectangle can have squared distance $\leq \tau$. Moreover, querying a promoted local R*-tree for its $k$ nearest entries suffices since its $(k+1)$-st nearest entry cannot belong to the global top-$k$.
\end{proof}

\section{Workload and Verification}

The primary workload comprises three deterministic 12,000-operation traces (\emph{moving-hotspot}, \emph{write-heavy}, and \emph{insert-only}) derived from 2024 U.S.\ Census TIGER/Line road geometries~\cite{census2024} for each of 24 counties ($72$ primary traces at capacity 32), plus four larger 60,000-operation write-heavy scaling traces at capacity 64 ($76$ traces total, $152$ trace-by-reinsertion evaluations). Operations cover insertion, deletion, motion, inclusive range intersection, and $k$-nearest-neighbor queries. Every method replays the identical trace text and verifies its SHA-256 digest. Within each county we sum the median replay times of the three profiles before forming the paired ratio against the baseline; build-inclusive times are reported separately, and 95\% bootstrap intervals resample counties as independent units.

At verification checkpoints, the harness audits every cell's live entries, ownership (\texttt{locate}), conservative bounds, descendant counts, and local R*-tree invariants. Every primary query (316,717 per reinsertion setting) is checked against an independent exhaustive scan over all active rectangles. On the four scaling traces, every query signature (83,690 per setting) is compared against its archived baseline and 2,265 sampled queries are verified by exhaustive scan.

\section{Results and Limitations}

In the preceding strict-directory study (capacity 32), median replay speedups over the pruned single R*-tree were $2.01\times$ ($1.94\text{--}2.09\times$) without reinsertion and $2.65\times$ ($2.44\text{--}2.82\times$) with 30\% reinsertion ($2.21\times$ and $3.01\times$ including initial build).

\begin{table}[ht]
\centering
\small
\caption{Capacity-32 paired median speedups of the loose R*-tree directory across 24 U.S.\ counties (95\% bootstrap intervals in brackets where aggregated across profiles).}
\label{tab:spatial}
\begin{tabularx}{\linewidth}{>{\raggedright\arraybackslash}Xcc}
\toprule
\textbf{Loose placement, capacity 32} & \textbf{Reinsertion off} & \textbf{30\% reinsertion} \\
\midrule
Replay vs.\ pruned single tree & $2.55\times$ [$2.44\text{--}2.68$] & $3.34\times$ [$3.11\text{--}3.46$] \\
Replay vs.\ strict forest (same run) & $1.26\times$ & $1.25\times$ \\
Replay and initial build vs.\ single tree & $2.77\times$ & $3.72\times$ \\
Moving hotspot vs.\ single tree & $1.72\times$ & $1.85\times$ \\
Write-heavy vs.\ single tree & $2.67\times$ & $3.17\times$ \\
Insert-only vs.\ single tree & $3.98\times$ & $6.48\times$ \\
\bottomrule
\end{tabularx}
\end{table}

All 152 paired evaluations matched the archived query signatures and passed independent verification. At capacity 64 across the four scaling traces, loose placement achieved median speedups of $2.72\times$ (reinsertion off) and $3.28\times$ (30\% reinsertion) over the single tree ($1.31\times$ and $1.51\times$ over the strict directory). Because loose placement routes by center and size rather than strict cell containment, a small rectangle crossing a grid line stays at its deep level instead of migrating to a high ancestor cell; while whole-replay timing does not isolate this mechanism alone, it aligns with the consistent $1.25\text{--}1.26\times$ advantage over the strict directory.

These measurements cover single-threaded, in-memory planar road rectangles whose coordinate magnitudes are checked before enabling specialized planar arithmetic, with a fixed grid domain drawn from initial insertions. Out-of-core storage, drifting geographic extents, higher dimensions, or custom nearest-neighbor tie-breaking contracts could alter both performance and semantics.

\paragraph{Code and data availability.}
The directory implementation (\nolinkurl{native/rtree_forest.hpp}), replay and verification harness (\nolinkurl{native/rtree.cpp}, \nolinkurl{native/rtree_policy.hpp}), tests (\nolinkurl{tests/test_native.py}), benchmark scripts (\nolinkurl{benchmarks/twice_spatial.py}), and archived journals (\nolinkurl{manuscripts/runs/spatial.jsonl}, \nolinkurl{manuscripts/summary.json}) are included in the repository.

\begin{thebibliography}{3}
\bibitem{beckmann1990}
N.~Beckmann, H.-P. Kriegel, R.~Schneider, and B.~Seeger.
\newblock The R*-tree: An efficient and robust access method for points and rectangles.
\newblock In \emph{Proceedings of the ACM SIGMOD International Conference on Management of Data}, pages 322--331, 1990.
\newblock \href{https://doi.org/10.1145/93597.98741}{doi:10.1145/93597.98741}.

\bibitem{samet2013}
H.~Samet, J.~Sankaranarayanan, and M.~Auerbach.
\newblock Indexing methods for moving object databases: Games and other applications.
\newblock In \emph{Proceedings of the ACM SIGMOD International Conference on Management of Data}, pages 169--180, 2013.
\newblock \href{https://doi.org/10.1145/2463676.2465332}{doi:10.1145/2463676.2465332}.

\bibitem{census2024}
U.S.\ Census Bureau.
\newblock 2024 TIGER/Line Shapefiles: Roads.
\newblock \href{https://www2.census.gov/geo/tiger/TIGER2024/ROADS/}{www2.census.gov/geo/tiger/TIGER2024/ROADS/}.
\end{thebibliography}

\end{document}
